\RequirePackage{fix-cm}
\documentclass[12pt]{article}
\usepackage[
  paperwidth=210mm,
  paperheight=248mm,
  left=20mm,
  right=6mm,
  top=5mm,
  bottom=5mm,
  includeheadfoot,
  headheight=16pt,
  headsep=3mm,
  footskip=8mm
]{geometry}


\usepackage{fancyhdr}
\pagestyle{fancy}
\fancyhf{}
\fancyfoot[C]{\thepage}
\renewcommand{\headrulewidth}{0pt}

\usepackage{hyperref}
\hypersetup{hidelinks}

\linespread{1.2}
\setlength{\emergencystretch}{3em}

\usepackage{amsmath,amssymb,booktabs,array}
\newcommand{\F}{\mathbb F_2}
\newcommand{\wt}{\operatorname{wt}}
\newcommand{\rk}{\operatorname{rank}}
\newcommand{\supp}{\operatorname{supp}}
\newcommand{\dd}{\mathrm d}
\title{Increasing Lifted-Product Code Distance\\by Reusing Fixed Constituent Expressions}
\author{Kenta Kasai\\\small Institute of Science Tokyo}
\date{November 24, 2026}
\renewcommand{\abstractname}{Abstract}
\renewcommand{\refname}{References}
\renewcommand{\tablename}{Table}

\begin{document}
\maketitle
\thispagestyle{fancy}

\begin{abstract}
Our previous work extended circulant- and affine-permutation pair-partition
(CPM--PP and APM--PP) codes by retaining the pair partition and fixed
integer parameters while changing the size. Here we pursue the same
objective for lifted-product (LP) codes: reuse the structure of an existing
code to obtain a longer code with greater distance. For LP codes over
finite group algebras, we retain the constituent matrix sizes and the
expressions defining their entries, including all terms, integer exponents
and multiplication order, and evaluate them in a larger finite group.
Commuting left and right regular actions preserve CSS orthogonality.
Distinct evaluated terms preserve check weights, while binary ranks
determine the actual rate.
As a concrete application, we extend the Mitten code from $[[975,195,24]]$
to $[[1650,330,30]]$ and $[[2500,500,30]]$ at unchanged rate 20\%
and row weight nine. The larger example also preserves the
square-invertibility condition, giving a complete canonical logical basis
generated by group actions.
The exact distances follow from exhaustive exclusion through weight 29
in both sectors and a verified weight-30 logical operator.
These examples establish finite-length distance improvement by reusing fixed
constituent expressions;
selecting evaluation groups for further improvement remains an open problem.
\end{abstract}

\section{Background and objective}
Our previous work on circulant- and affine-permutation pair-partition
(CPM--PP and APM--PP) codes \cite[Secs.~2--3]{explicit} was motivated by
a design question: can we obtain longer codes with greater minimum
distance by retaining much of an existing code's structure and changing
only selected parameters? The aim was to reuse the structure of a good
code while avoiding parameter choices that produce low-weight nontrivial
logical operators. We retained the pair partition and fixed integer
parameters extracted from an original code, and regenerated only the
size-dependent arrays at a new size. This preserved the orthogonality
mechanism and check weights while providing candidates for distance
improvement.

Quantum LDPC codes limit both the number of qubits in a check and the
number of checks involving a qubit \cite{sparse}. In the CSS construction,
the X and Z checks must commute \cite{css,steane}, while low-weight
nontrivial logical operators must be avoided. Reducing the physical-qubit cost of reliable logical information requires
both a substantial encoding rate and protection against increasingly
large errors. Quantum LDPC constructions have transformed the attainable
rate--distance tradeoff: hypergraph products achieve positive asymptotic
rate and square-root distance \cite{hgp}, while LP codes over suitable
non-Abelian groups achieve both dimension and distance proportional to
blocklength \cite{good}. Quantum Tanner codes provide a related route to
asymptotically good codes \cite{tanner}; their recent generalization to
non-Cayley Schreier graphs enlarges the available design space
\cite{generalizedtanner}. These existence results establish what is possible
as blocklength grows. Finding useful, explicitly specified finite codes
with simultaneously high rate, large distance and sparse checks remains
a separate design problem.

At finite length, generalized hypergraph products give structured
quasi-cyclic examples \cite[Sec.~5]{finite}, and lifted products separate
small constituent matrices from the algebra used to expand them into
binary checks \cite[Sec.~III]{lp}. Balanced products \cite{balanced} and
two-block group-algebra codes \cite{twoblock} further develop this use of
algebraic structure. Recent non-Abelian ZSZ-LP codes target design rate
$1/5$ with check weight nine; one reported instance has parameters
$[[625,125,\leq 20]]$, with the stated distance being an upper bound
\cite[Table~8]{zsz}. Structured-concept-evolution searches report the LP
codes $[[585,121,24]]$ at check weight 12 and $[[585,117,20]]$ at check
weight nine, with both distances certified exactly
\cite[Sec.~II and Supplemental Material S4--S5]{sce}.
The Mitten construction likewise exploits non-Abelian symmetries for
high-rate codes, and studies decoding, syndrome extraction and logical
operations \cite{mitten}. Its commuting left and right regular actions
ensure CSS orthogonality \cite[Appendix A.B--C]{mitten}.

Progress also extends beyond these LP examples. Bivariate bicycle codes
combine sparse checks with low-overhead memory protocols; the
$[[144,12,12]]$ example illustrates this finite-length direction
\cite{bicycle}. Its rate $1/12$ differs from the $1/5$ rate considered here.
High-rate hypergraph-product proposals for static neutral-atom registers
also connect code structure to realizable long-range interactions
\cite{neutral}. Such studies make the finite-length question especially
consequential: increasing distance is valuable only in the context of the
rate, blocklength, check weights and intended implementation.
Code distance and circuit-level logical error rate are distinct criteria;
recent LP searches explicitly find that optimizing them can favor
different codes \cite{sce}.

Our earlier APM construction controls commutativity to reconcile sparse,
regular checks with CSS orthogonality \cite{barrier}. PP constructions
instead pair block-column contributions so that they cancel
\cite{pp}; subsequent work develops the finite-length tradeoffs among
column weight, check weight, rate and distance \cite{design}.
Complete low-weight canonical logical bases have also been constructed
for a quaternary extension of PP codes \cite{ppbasis}. This line of work
motivates retaining useful construction information when increasing size,
while checking which distance and logical-basis properties survive.

Larger size alone need not remove structural distance restrictions.
Classical quasi-cyclic LDPC codes have block-structure-dependent distance
upper bounds \cite{qcbound}; for full single-CPM arrays, a fixed integer
exponent matrix can attain the structural bound at every sufficiently
large lift size \cite{cpmattain}. These are classical results, and their
distance conclusions do not transfer directly to quantum LP codes.
The sampling rule matters as well: for full permutation-block CSS arrays
sampled uniformly subject to orthogonality, relative distance vanishes
in probability at fixed column weight $J\geq 3$ and even row weight
$L>2J$ as the permutation size grows \cite{randomcss}.
That result concerns the specified ensemble; it does not establish a
distance limit for structured LP constructions with sums and zero blocks.
The question here is whether successful constituent expressions can be
reused at larger sizes while preserving their algebraic advantages.

We apply this strategy to LP codes by retaining the constituent matrix
sizes and the ordered expressions defining their entries while changing
the finite evaluation group. We assess the resulting distance together
with rate, blocklength and check weights \cite{design}.

Evaluating the fixed Mitten constituent expressions at $G(25,20,4)$ gives
$[[2500,500,30]]$, increasing the original distance from 24 to 30 while
retaining rate $1/5$ and check-row weight nine. The two required
constituent blocks are invertible, so this evaluation also retains the
Mitten canonical logical basis generated by group actions
\cite[Definition~4 and Theorem~1]{mitten}.
The exact distance is established by exhaustive exclusion through weight 29
and a verified Z logical operator of weight 30.

Enlarging an LP code while retaining structure is an established direction.
Extensions by group extensions and graph lifts preserve local Tanner-graph
structure, relate logical operators through chain and cochain maps, and
provide finite codes with improved distance \cite{lifting}.
That construction uses a surjection from the target group onto the parent
\cite[Sec.~II.1]{lifting}. The Mitten design pipeline also already varies
the group and constituent matrices to seek good codes \cite{mitten}.
We position the present contribution as the reuse of a successful set of
constituent expressions under a restricted change of evaluation group,
with explicit checks of the retained algebraic properties and exact
verification of the resulting distance improvement. Every term, integer
exponent and multiplication order is fixed.

This reuse permits evaluations outside the group-extension relation above.
For the original Mitten code and the $[[2500,500,30]]$ example, the group
orders are 195 and 500, respectively. Since 195 does not divide 500,
there is no surjective homomorphism from the target group onto the parent.
Consequently, this evaluation is not a group-extension lift of that parent
in the cited framework. Distance improvement is verified for the selected evaluations; identifying
which evaluation groups improve distance remains a design problem.

For LP codes, the structure that ensures orthogonality differs from the
pair partition used in the previous CPM--PP and APM--PP extensions.
The Mitten checks contain zero blocks and sums of permutation matrices,
so the previous method's single-permutation-block assumption does not
cover this example. Here the commuting left and right regular actions
provide the orthogonality mechanism, allowing the fixed constituent
expressions to be reused at a new group evaluation.

Distance verification is essential to this program. Classical approaches
include probabilistic low-weight searches \cite{leon,stern} and the
Brouwer--Zimmermann method using multiple information sets
\cite{grassl,lisonek}, with distributed-memory implementations also
available \cite{parallelbz}. For quantum codes, the search must distinguish
stabilizers from nontrivial logical operators. Relevant approaches include
random-information-set upper-bound searches \cite{qdistrnd}, fast
Brouwer--Zimmermann implementations for symplectic weight \cite{quantumbz},
and connected-cluster enumeration exploiting sparse checks \cite{clusters}.
Recent benchmarks also compare exact SAT and mixed-integer methods
\cite{distancebench}. Our work on APM distance witnesses emphasizes that
an upper bound requires an explicit kernel vector outside the stabilizer
row space \cite{apmwitness}; exactness additionally requires excluding
all lighter nontrivial logical operators.

We describe the method for arbitrary constituent matrix sizes and finite
evaluation groups, and use the fixed Mitten expressions as a worked example.
Distance improvement is established by exhaustive exclusion using
connected-cluster search \cite{clusters} and LP-specific reductions
\cite[Sec.~3.6]{qedist}, together with a logical operator attaining the
lower bound. The PP construction and its later design framework also give
exhaustive distance-verification procedures \cite{pp,design}.
Section~\ref{sec:construction} gives the common construction;
Section~\ref{sec:results} specifies the example and verifies its distance.

\section{Extending LP codes with fixed constituent expressions}\label{sec:construction}
Reusing an LP design requires specifying which data remain fixed and
which properties follow from those data. For arbitrary constituent matrix
sizes over finite group algebras, commuting regular actions preserve CSS
orthogonality; term distinctness and binary ranks determine the retained
weights and rate. Section~\ref{sec:results} applies these conditions to
the Mitten constituent expressions.

Let $\alpha$ and $\beta$ have sizes $r_A\times s_A$ and
$r_B\times s_B$, respectively. Their entries are fixed binary sums of
products of named generators and their inverses. Fix every term, integer
exponent and factor order. An evaluation specifies a finite group $G$ and
the group elements assigned to those generators. Evaluating the same
expressions in the parent group and a larger target group gives the two
codes. The matrix sizes and expressions remain fixed, while the evaluated
group elements and their regular permutation matrices can change.
The chosen formal expressions are part of the input: expressions that
agree in the parent group can differ in the target group. We therefore
retain their displayed representatives as construction data.

Write $q=|G|$. In the binary group algebra $\F[G]$, coefficients add
modulo two and multiplication extends the group product bilinearly.
Collecting the coefficient of each $g,h\in G$ gives binary matrices
$\alpha_g$ of size $r_A\times s_A$ and $\beta_h$ of size
$r_B\times s_B$, so that $\alpha=\sum_g\alpha_g g$ and
$\beta=\sum_h\beta_h h$. For standard coordinate vectors $\mathbf e_u$,
$u\in G$, the regular permutation matrices are
\begin{equation}
 L_g\mathbf e_u=\mathbf e_{gu},\qquad R_h\mathbf e_u=\mathbf e_{uh}.
 \label{eq:regular}
\end{equation}
The binary constituent matrices are
\begin{equation}
 A=\sum_{g\in G}\alpha_g\otimes L_g,\qquad
 B=\sum_{h\in G}\beta_h\otimes R_h.
 \label{eq:constituents}
\end{equation}
Here $\otimes$ is the ordinary Kronecker product of binary matrices.
The LP check matrices \cite[Sec.~III]{lp} are
\begin{align}
 H_X&=\left[\,
 \sum_{g\in G}\alpha_g\otimes I_{s_B}\otimes L_g
 \;\middle|\;
 \sum_{h\in G}I_{r_A}\otimes\beta_h^T\otimes R_h^T
 \,\right],\label{eq:hx}\\
 H_Z&=\left[\,
 \sum_{h\in G}I_{s_A}\otimes\beta_h\otimes R_h
 \;\middle|\;
 \sum_{g\in G}\alpha_g^T\otimes I_{r_B}\otimes L_g^T
 \,\right].\label{eq:hz}
\end{align}
The identity factors repeat a constituent block across the index of the
other constituent. The two qubit blocks contain $s_As_Bq$ and $r_Ar_Bq$
qubits; the check-row counts are $r_As_Bq$ and $s_Ar_Bq$.

CSS orthogonality is preserved by every evaluation. Associativity gives
$L_gR_h=R_hL_g$, and $R_h^T=R_{h^{-1}}$, so each pair of terms in
$H_XH_Z^T$ contributes
$\alpha_g\otimes\beta_h^T\otimes
(L_gR_h^T+R_h^TL_g)=0$ over $\F$.
This uses the commuting regular actions also employed in the Mitten
construction \cite[Appendix A.B--C]{mitten}.
No homomorphism between the two evaluation groups is needed.

The reuse of fixed expressions also controls sparsity. If the terms of
each entry evaluate to distinct group elements in both groups, each
constituent block retains its row and column weight. Consequently, the
row and column weight patterns of $H_X,H_Z$ are retained.
The length and dimension are
\begin{equation}
 n=(s_As_B+r_Ar_B)q,\qquad
 k=n-\rk H_X-\rk H_Z.\label{eq:rate}
\end{equation}
The actual rate depends on the binary ranks. When both check matrices
have full row rank at both evaluations, the unchanged constituent sizes
give the same rate.

Changing the group alters the syndrome and stabilizer conditions on small
supports and can remove the parent's low-weight logical operators.
The construction guarantees orthogonality; distance improvement is
established on the resulting matrices by excluding all lighter nontrivial
logical operators. The following example carries out this verification.

\section{Mitten example and distance verification}\label{sec:results}
We apply the fixed-expression method with $r_A=r_B=1$ and $s_A=s_B=2$.
The two constituent rows are those of the original Mitten code, and only
its evaluation group changes.

The original code uses $G(13,15,9)$, and the two extensions use
$G(22,15,3)$ and $G(25,20,4)$. Write
\begin{equation}
 G(p,m,\rho)=\langle x,w\mid x^p=w^m=1,\ wxw^{-1}=x^\rho\rangle,
 \qquad (p,m,\rho)\in\{(13,15,9),(22,15,3),(25,20,4)\}.\label{eq:group}
\end{equation}
All three choices satisfy $\gcd(\rho,p)=1$ and $\rho^m\equiv1\pmod p$,
so conjugation by $w$ defines an invertible action consistent with
$w^m=1$. The group has order $q=pm$, and each element has unique
coordinates $(i,j)=x^iw^j$, with $0\leq i<p$ and $0\leq j<m$.
Multiplication and inversion are given by
\begin{equation}
 (i,j)(u,v)=(i+\rho^ju\bmod p,j+v\bmod m),\qquad
 (i,j)^{-1}=(-\rho^{-j}i\bmod p,-j\bmod m).
 \label{eq:groupcalc}
\end{equation}
We order these coordinates lexicographically by $(i,j)$.

The constituent rows are $\alpha=(\alpha_{00},\alpha_{01})$ and
$\beta=(\beta_{00},\beta_{01})$. In the coordinates of
Eq.~\eqref{eq:groupcalc}, the entries of the original code
\cite[Appendix K, Table XIII]{mitten} are
% Generated by generate_construction_table.py; do not edit.
\begin{equation}
\begin{aligned}
\alpha_{00}&=x^{11}w^{7}+x^{11}w^{11}+x^{5}w^{14},\\
\alpha_{01}&=1+x^{9}+x^{4}w^{14},\\
\beta_{00}&=w^{-7}x^{-1}+w^{-3}x^{-3}+w^{-12}x^{-9},\\
\beta_{01}&=1+w^{-2}x^{-9}+w^{-14}x^{-6}.
\end{aligned}
\label{eq:mittenentries}
\end{equation}

The same terms, integer exponents and factor order are used for all three
groups. In particular, the exponents are retained as written, rather than
replaced by congruent representatives for the original group.
Changing $(p,\rho)$ from $(13,9)$ to $(22,3)$ keeps the order of $w$ at
15 and changes the order of its conjugation action from three to five.
The group orders are 195 and 330; since $195\nmid330$, the latter group
does not map onto the former.

The inverses in $\beta$ match the positive right-action convention in
Eq.~\eqref{eq:regular}. For these constituent sizes,
Eqs.~\eqref{eq:hx}--\eqref{eq:hz} give two qubit blocks of sizes $4q$ and
$q$. The first has coordinates $(a,b,u)$, with $a,b\in\{0,1\}$ and
$u\in G$, in that order; the second has coordinates $u\in G$.
The left part of $H_X$ has block rows $[A_{00}\ 0\ A_{01}\ 0]$ and
$[0\ A_{00}\ 0\ A_{01}]$, and its right part stacks $B_{00}^T,B_{01}^T$.
At the original group parameters, these matrices reproduce the supplied
Mitten checks under the common qubit and check-row permutations recorded
in \cite{certificates}.

In all three groups the terms of every constituent entry remain distinct.
Each constituent block has row and column weight three, giving check-row
weight nine and column weights three and six in the first and second
qubit blocks. Binary Gaussian elimination verifies full row rank for both
checks of each code, so Eq.~\eqref{eq:rate} gives $k=q$ and rate $1/5$.

Table~\ref{tab:comparison} lists the
original Mitten code and its two extensions. The original distance is certified
in \cite[Appendix C]{qedist}; we reproduce its matrices and verify a
logical operator attaining that distance.
\begin{table}[htbp]
\centering
\small
\caption{The original Mitten code and two fixed-expression extensions, evaluated over $G(p,m,\rho)$ in Eq.~\eqref{eq:group} using the fixed entries in Eq.~\eqref{eq:mittenentries}. Column and row weight list the values in each check matrix; $g$ is their common Tanner girth. All distances are exact. The original distance is from \cite{qedist}. Code data for all three rows and complete verification records for both extensions are in \cite{certificates}. The 2500-qubit extension retains the defining Mitten invertibility condition; the 1650-qubit extension does not.}
\label{tab:comparison}
\begin{tabular}{lrrrrrrr}
\toprule
Construction&$p$&$m$&$\rho$&$[[n,k,d]]$&Column weight&Row weight&$g$\\
\midrule
Mitten LP (original code)~\cite{mitten}&13&15&9&$[[975,195,24]]$&$\{3,6\}$&$9$&$4$\\
Mitten LP (extension)&22&15&3&$[[1650,330,30]]$&$\{3,6\}$&$9$&$4$\\
Mitten LP (extension)&25&20&4&$[[2500,500,30]]$&$\{3,6\}$&$9$&$4$\\
\bottomrule
\end{tabular}
\end{table}


For the original code and the $[[2500,500,30]]$ extension, the relative
distances are approximately $2.46\%$ and $1.20\%$, respectively.
The absolute distance increases by $25\%$ as the blocklength increases
from 975 to 2500. Both codes use five data qubits per logical qubit;
this rate-based ratio does not include syndrome ancillas or logical-operation
gadgets. The gain in absolute distance thus comes with increased
blocklength and reduced relative distance. Decoding and circuit-level
performance require separate evaluation.

For $G(22,15,3)$, the exhaustive computation described below excludes nontrivial logical
operators through weight 29 in both sectors. A verified X logical operator
has weight 30, establishing the exact distance in
Table~\ref{tab:comparison}. More specifically, $d_X=30$ and
$30\leq d_Z\leq38$. The global distance is $\min(d_X,d_Z)=30$.

The evaluation at $G(25,20,4)$ also gives exact distance 30.
Both check matrices have rank 1000, and exhaustive exclusion through
weight 29 in both sectors is matched by a verified weight-30 Z logical
operator. Thus $d_Z=30$ and $30\leq d_X\leq122$.
This evaluation additionally retains the square-invertibility condition
for Mitten codes: $A_{01}$ and $B_{01}$ have rank 500
\cite[Definition~4 and Theorem~1]{mitten}.
The resulting 500 canonical logical pairs are generated by the group
actions from X and Z seeds of weights 256 and 262, respectively.
The $G(22,15,3)$ evaluation has ranks 320 for these two blocks of size
330 and therefore requires the direct rank and distance checks above.

These seed weights describe the constructed canonical basis, whereas
minimum distance is minimized over all nontrivial logical operators.
Finding a complete canonical basis with lighter representatives remains
an optimization problem.

The group relations illustrate how fixed expressions can behave differently
at the new evaluation. In $G(13,15,9)$, $9^3\equiv1\pmod{13}$, so
$w^3xw^{-3}=x$. In $G(25,20,4)$, $4^3\equiv14\pmod{25}$, so
$w^3xw^{-3}=x^{14}\ne x$. Thus a relation available to close a word in
the parent group is lost. This changes possible cancellation patterns in
syndrome equations. A direct identification of which parent weight-24
logical supports are obstructed by these changed relations remains to
be established; the exact distance statement rests on the exhaustive
verification below.

Candidate selection and distance certification serve different roles.
The constituent expressions and their integer exponents are fixed; the
evaluation parameters $(p,m,\rho)$ are varied. Rank, term-distinctness
and invertibility tests screen the algebraic properties of a candidate.
We report selected evaluations with verified algebraic properties and
completed distance verification. Finding a logical
operator supplies an upper bound; an unsuccessful heuristic search or
an unfinished exhaustive run supplies no certified lower bound.

To certify a distance lower bound, we exclude all lighter nontrivial logical operators.
For the Z sector the search space is $\ker H_X\setminus\operatorname{row}H_Z$;
X and Z are interchanged for the other sector.
Each logical operator providing an upper bound is checked for zero syndrome
and nonmembership in the stabilizer row space.

We use connected-cluster distance search \cite{clusters} and the LP-specific
reduction of QEDist \cite[Sec.~3.6]{qedist}.
Connect two qubits when they occur in a common
syndrome check. A zero-syndrome support decomposes into connected components
in this graph, each with zero syndrome. If the operator is nontrivial, at least
one component is nontrivial and has no greater weight.
For the LP reduction, write a Z operator as $(\mathbf u,\mathbf v)$
in the two qubit blocks of Eqs.~\eqref{eq:hx}--\eqref{eq:hz}.
If $\mathbf v\ne0$, the full-code search starts from root qubits in the second block.
Automorphisms preserving the sector reduce the number of root qubits:
every nonempty support
in this block can be moved to contain an orbit representative. Excluding
earlier orbits partitions the search cases.

If $\mathbf v=0$, the syndrome condition separates into constituent columns
$\mathbf a_b\in K=\ker A$. Let $E_b$ embed such a column in the first block.
Since these columns have disjoint supports, a nontrivial
$\mathbf u=\sum_bE_b\mathbf a_b$ has a nontrivial summand of no greater weight.
Define

\begin{equation}
 S_{\cap}=\{\mathbf a\in K:E_b\mathbf a\in\operatorname{row}H_Z
                 \quad\forall b\}.\label{eq:commontrivial}
\end{equation}

Searching $K\setminus S_{\cap}$ covers every such summand. For each
resulting vector $\mathbf a$, verify an embedding $E_b\mathbf a$ outside
$\operatorname{row}H_Z$. Both searches are required. The X-sector argument
uses the other constituent matrix. The software checks the LP block
structure of the binary matrices, the subspaces and the code automorphisms.

We check the search partitions and completed execution records against
the implementation, input matrices and code automorphisms. We also verify
that the refined residue classes are disjoint and cover the search space.
Only completed exhaustive searches contribute to the reported lower bounds;
their records identify the source version and input matrices.

The code specifications, check matrices, logical operators and
reconstruction procedures for the original code and both extensions,
together with the complete distance-verification records for both
extensions, are available in the companion repository~\cite{certificates}.

\section{Conclusion and future directions}
Continuing the structure-preserving extension approach previously used
for CPM--PP and APM--PP codes, we have considered group-algebra LP codes.
Here fixed constituent expressions can be reused while preserving CSS
orthogonality through commuting regular actions. The operation applies across
constituent matrix sizes and evaluation groups. Its application to the
Mitten code increases the exact distance at unchanged actual rate and
check weights, as established by exhaustive computation. The
$[[2500,500,30]]$ evaluation also retains the required constituent
invertibility and complete group-generated canonical logical basis.

The remaining question is which group parameters make this change effective.
Understanding how the original low-weight logical operators change could
guide parameter selection and reduce verification cost. Comparisons with
other constructions at matched length, rate and check weights, and the
extension's decoder performance, remain to be investigated.

\begingroup
\footnotesize
\begin{thebibliography}{99}
\interlinepenalty=10000
\setlength{\itemsep}{0pt}
\setlength{\parsep}{0pt}
\bibitem{explicit}K. Kasai, ``Increasing Quantum LDPC Code Distance through
Explicit Extensions,'' arXiv:2611.28847v1 (2026),
\url{https://arxiv.org/abs/2611.28847v1}.
\bibitem{sparse}D. J. C. MacKay, G. Mitchison, and P. L. McFadden,
``Sparse-graph codes for quantum error correction,'' \emph{IEEE Transactions
on Information Theory}, vol.~50, no.~10, pp.~2315--2330, 2004,
\url{https://doi.org/10.1109/TIT.2004.834737}.
\bibitem{css}A. R. Calderbank and P. W. Shor, ``Good quantum error-correcting
codes exist,'' \emph{Physical Review A}, vol.~54, pp.~1098--1105, 1996,
\url{https://doi.org/10.1103/PhysRevA.54.1098}.
\bibitem{steane}A. M. Steane, ``Error correcting codes in quantum theory,''
\emph{Physical Review Letters}, vol.~77, pp.~793--797, 1996,
\url{https://doi.org/10.1103/PhysRevLett.77.793}.
\bibitem{hgp}J.-P. Tillich and G. Z\'emor,
``Quantum LDPC Codes With Positive Rate and Minimum Distance Proportional
to the Square Root of the Blocklength,'' \emph{IEEE Transactions on Information Theory},
vol.~60, no.~2, pp.~1193--1202, 2014,
\url{https://doi.org/10.1109/TIT.2013.2292061}.
\bibitem{good}P. Panteleev and G. Kalachev,
``Asymptotically Good Quantum and Locally Testable Classical LDPC Codes,''
in \emph{Proceedings of the 54th Annual ACM SIGACT Symposium on Theory of Computing},
pp.~375--388, 2022, \url{https://doi.org/10.1145/3519935.3520017}.
\bibitem{tanner}A. Leverrier and G. Z\'emor, ``Quantum Tanner codes,''
in \emph{2022 IEEE 63rd Annual Symposium on Foundations of Computer Science},
pp.~872--883, 2022, \url{https://doi.org/10.1109/FOCS54457.2022.00117}.
\bibitem{generalizedtanner}O. {\AA}. Mostad, E. Rosnes, and H.-Y. Lin,
``Asymptotically Good Generalized Quantum Tanner Codes,''
\emph{IEEE Journal on Selected Areas in Information Theory}, vol.~6,
pp.~367--382, 2025, \url{https://doi.org/10.1109/JSAIT.2025.3594310}.
\bibitem{finite}P. Panteleev and G. Kalachev,
``Degenerate Quantum LDPC Codes With Good Finite Length Performance,''
\emph{Quantum}, vol.~5, article~585, 2021,
\url{https://doi.org/10.22331/q-2021-11-22-585}.
\bibitem{lp}P. Panteleev and G. Kalachev, ``Quantum LDPC Codes With Almost
Linear Minimum Distance,'' \emph{IEEE Transactions on Information Theory},
vol.~68, no.~1, pp.~213--229, 2022,
\url{https://doi.org/10.1109/TIT.2021.3119384}.
\bibitem{balanced}N. P. Breuckmann and J. N. Eberhardt,
``Balanced Product Quantum Codes,'' \emph{IEEE Transactions on Information Theory},
vol.~67, no.~10, pp.~6653--6674, 2021,
\url{https://doi.org/10.1109/TIT.2021.3097347}.
\bibitem{twoblock}H.-K. Lin and L. P. Pryadko,
``Quantum two-block group algebra codes,'' \emph{Physical Review A},
vol.~109, no.~2, article~022407, 2024,
\url{https://doi.org/10.1103/PhysRevA.109.022407}.
\bibitem{zsz}Y. Hong, ``Quantum LDPC codes with design rate 1/5 and good
performance below 1000 physical qubits,'' arXiv:2607.27644v1 (2026),
\url{https://arxiv.org/abs/2607.27644v1}.
\bibitem{sce}Z. Liu and F. Marquardt, ``Large-Language-Model Discovery of
Quantum LDPC Codes through Structured Concept Evolution,''
arXiv:2606.24808v2 (2026), \url{https://arxiv.org/abs/2606.24808v2}.
\bibitem{mitten}A. Bhardwaj, M. Ma, N. Meister, R. King, D. Bluvstein,
J. Preskill, M. Cain, Q. Xu, and H.-Y. Huang,
``High-rate qLDPC processors,'' arXiv:2607.28795v2 (2026),
\url{https://arxiv.org/abs/2607.28795v2}.
\bibitem{bicycle}S. Bravyi, A. W. Cross, J. M. Gambetta, D. Maslov,
P. Rall, and T. J. Yoder, ``High-threshold and low-overhead fault-tolerant
quantum memory,'' \emph{Nature}, vol.~627, pp.~778--782, 2024,
\url{https://doi.org/10.1038/s41586-024-07107-7}.
\bibitem{neutral}L. Pecorari, S. Jandura, G. K. Brennen, and G. Pupillo,
``High-rate quantum LDPC codes for long-range-connected neutral atom
registers,'' \emph{Nature Communications}, vol.~16, article~1111, 2025,
\url{https://doi.org/10.1038/s41467-025-56255-5}.
\bibitem{barrier}K. Kasai, ``Breaking the Orthogonality Barrier in Quantum
LDPC Codes,'' \emph{Quantum}, vol.~10, article~2205, 2026,
\url{https://doi.org/10.22331/q-2026-09-09-2205}.
\bibitem{pp}K. Okada and K. Kasai, ``Pair-Partition Constructions for
CPM-Based Quantum LDPC Codes,'' arXiv:2607.14091v4 (2026),
\url{https://arxiv.org/abs/2607.14091v4}.
\bibitem{design}J. Y. Lee, K. Okada, N. Maskara, K. Kasai, and H. Zhou,
``Design Principles for Ultra-High-Rate Quantum Codes,''
arXiv:2609.30069v1 (2026), \url{https://arxiv.org/abs/2609.30069v1}.
\bibitem{ppbasis}K. Okada and K. Kasai, ``Low-Weight Canonical Logical
Bases from Pair-Partition Codes,'' arXiv:2609.35601v1 (2026),
\url{https://arxiv.org/abs/2609.35601v1}.
\bibitem{qcbound}R. Smarandache and P. O. Vontobel, ``Quasi-cyclic LDPC
codes: Influence of proto- and Tanner-graph structure on minimum Hamming
distance upper bounds,'' \emph{IEEE Transactions on Information Theory},
vol.~58, no.~2, pp.~585--607, 2012,
\url{https://doi.org/10.1109/TIT.2011.2173244}.
\bibitem{cpmattain}K. Kasai, ``CPM-LDPC Codes Attaining the Minimum-Distance
Bound,'' arXiv:2609.16836v1 (2026),
\url{https://arxiv.org/abs/2609.16836v1}.
\bibitem{randomcss}K. Kasai, ``Vanishing Relative Distance of Uniformly
Random CSS LDPC Codes Built from Permutation Matrices,''
arXiv:2611.21389v1 (2026), \url{https://arxiv.org/abs/2611.21389v1}.
\bibitem{lifting}Y. Hirasaki and J. Y. Lee, ``Lifting Lifted Product Codes,''
arXiv:2607.28621v1 (2026),
\url{https://arxiv.org/abs/2607.28621v1}.
\bibitem{leon}J. S. Leon, ``A probabilistic algorithm for computing minimum
weights of large error-correcting codes,'' \emph{IEEE Transactions on
Information Theory}, vol.~34, no.~5, pp.~1354--1359, 1988,
\url{https://doi.org/10.1109/18.21270}.
\bibitem{stern}J. Stern, ``A method for finding codewords of small weight,''
in \emph{Coding Theory and Applications}, Lecture Notes in Computer Science,
vol.~388, pp.~106--113, Springer, 1989,
\url{https://doi.org/10.1007/BFb0019850}.
\bibitem{grassl}M. Grassl, ``Searching for linear codes with large minimum
distance,'' in \emph{Discovering Mathematics with Magma}, Algorithms and
Computation in Mathematics, vol.~19, pp.~287--313, Springer, 2006,
\url{https://doi.org/10.1007/978-3-540-37634-7_13}.
\bibitem{lisonek}P. Lison\v{e}k and L. Trummer, ``Algorithms for the minimum
weight of linear codes,'' \emph{Advances in Mathematics of Communications},
vol.~10, no.~1, pp.~195--207, 2016,
\url{https://doi.org/10.3934/amc.2016.10.195}.
\bibitem{parallelbz}G. Quintana-Ort\'{i}, F. Hernando, and F. D. Igual,
``Algorithm 1033: Parallel implementations for computing the minimum
distance of a random linear code on distributed-memory architectures,''
\emph{ACM Transactions on Mathematical Software}, vol.~49, no.~1,
pp.~1--24, 2023, \url{https://doi.org/10.1145/3573383}.
\bibitem{qdistrnd}L. P. Pryadko, V. A. Shabashov, and V. K. Kozin,
``QDistRnd: A GAP package for computing the distance of quantum
error-correcting codes,'' \emph{Journal of Open Source Software}, vol.~7,
no.~71, article~4120, 2022, \url{https://doi.org/10.21105/joss.04120}.
\bibitem{quantumbz}F. Hernando, G. Quintana-Ort\'{i}, and M. Grassl,
``Fast Algorithms and Implementations for Computing the Minimum Distance
of Quantum Codes,'' \emph{ACM Transactions on Quantum Computing}, vol.~7,
no.~2, article~12, 2026, \url{https://doi.org/10.1145/3795877}.
\bibitem{clusters}I. Dumer, A. A. Kovalev, and L. P. Pryadko,
``Distance Verification for Classical and Quantum LDPC Codes,''
\emph{IEEE Transactions on Information Theory}, vol.~63, no.~7,
pp.~4675--4686, 2017,
\url{https://doi.org/10.1109/TIT.2017.2690381}.
\bibitem{distancebench}M. Webster, A. Jacob, and O. Higgott,
``Distance-Finding Algorithms for Quantum Codes and Circuits,''
arXiv:2603.22532v2 (2026), \url{https://arxiv.org/abs/2603.22532v2}.
\bibitem{apmwitness}K. Kasai, ``Heuristic Search for Minimum-Distance
Upper-Bound Witnesses in Quantum APM-LDPC Codes,'' arXiv:2604.15307v1 (2026),
\url{https://arxiv.org/abs/2604.15307v1}.
\bibitem{qedist}J. Y. Lee, ``QEDist: Practical Exact Distance Certification
for Stabilizer Codes,'' arXiv:2611.21709v1 (2026),
\url{https://arxiv.org/abs/2611.21709v1}.
\bibitem{certificates}K. Kasai, ``Fixed-word lifted-product code distance
certificates,'' version 1.0.0, 2026,
\url{https://github.com/kasaikenta/lp-fixed-word-certificates/releases/tag/v1.0.0}.
\end{thebibliography}
\endgroup

\end{document}
